% Source: https://artofproblemsolving.com/community/c3h2431035_500th_celebration_post
% Stable identity: aops-topic-2431035-q35
Let $s=\sqrt3+1$ and take $A=(0,0)$, $B=(s,0)$, $C=(s,s)$, $D=(0,s)$ (counterclockwise; the clockwise labelling gives the mirror image of the same region, with the same area). Then $P=(1,1)$, and
\[PA=\sqrt2,\qquad PB=PD=\sqrt{3+1}=2,\qquad PC=\sqrt2\cdot\sqrt3=\sqrt6.\]
Measure directions from $P$ counterclockwise from the positive $x$-direction. Then $B$, $C$, $D$, $A$ lie in the directions $-30^\circ$, $45^\circ$, $120^\circ$, $225^\circ$ (for instance $B-P=\left(\sqrt3,-1\right)$). The distances from $P$ to the lines $AB$ and $AD$ are $1$, and to $BC$ and $CD$ they are $\sqrt3$.

\textbf{How far the region reaches.} Let $r(\psi)$ be the distance from $P$ to the boundary of the square in direction $\psi$. The square is convex and contains $P$, so it meets each ray from $P$ in a segment of length $r(\psi)$. A point at distance $\rho$ from $P$ in direction $\varphi$ is swept exactly when, for some angle $\theta\in[0^\circ,90^\circ]$, the point at distance $\rho$ in direction $\varphi-\theta$ lies in the original square. So the swept region meets the ray in direction $\varphi$ in a segment of length
\[R(\varphi)=\max\{r(\psi):\ \varphi-90^\circ\le\psi\le\varphi\}.\]
Along a side of the square the distance from $P$ decreases to the foot of the perpendicular and then increases, so on any interval of directions the maximum of $r$ is attained at a vertex direction inside the interval or at an end of the interval. Going round $P$, this gives (the first column lists $\varphi$, and $r(\varphi)$ belongs to the starting square, $r(\varphi-90^\circ)$ to the final square $A'B'C'D'$):
\begin{itemize}
\item $45^\circ\le\varphi\le135^\circ$: $C$ is in the window, and $PC=\sqrt6$ is the largest distance of all, so $R=\sqrt6$: a quarter-circle traced by $C$.
\item $135^\circ\le\varphi\le150^\circ$: the only vertex in the window is $D$ ($PD=2$), while $r(\varphi-90^\circ)=\dfrac{\sqrt3}{\sin(\varphi-90^\circ)}\ge2$ (the ray hits side $CD$ between $C$ and $(2,s)$). So $R$ follows a side of the final square.
\item $150^\circ\le\varphi\le210^\circ$: $D$ is in the window; both end values are at most $2$, so $R=2$: an arc traced by $D$ (angle $60^\circ$).
\item $210^\circ\le\varphi\le225^\circ$: no vertex in the window; the larger end value is $r(\varphi-90^\circ)$, on side $DA$ between $D$ and $(0,2)$, so $R$ follows a side of the final square.
\item $225^\circ\le\varphi\le315^\circ$: $A$ is in the window and both end values are at most $\sqrt2$, so $R=\sqrt2$: a quarter-circle traced by $A$.
\item $315^\circ\le\varphi\le330^\circ$: no vertex in the window; the larger end value is $r(\varphi)$, on side $AB$ between $(2,0)$ and $B$: a side of the starting square.
\item $330^\circ\le\varphi\le390^\circ$: $B$ is in the window and both end values are at most $2$, so $R=2$: an arc traced by $B$ (angle $60^\circ$).
\item $30^\circ\le\varphi\le45^\circ$: $B$ is in the window, but $r(\varphi)=\dfrac{\sqrt3}{\cos\varphi}\ge2$ on side $BC$ between $(s,2)$ and $C$: a side of the starting square.
\end{itemize}

\begin{asy}
usepackage("amsmath");
size(7cm);
defaultpen(fontsize(10pt));
real s=sqrt(3)+1; pair P=(1,1);
pair A=(0,0), B=(s,0), C=(s,s), D=(0,s);
pair A2=(2,0), B2=(2,s), C2=(2-s,s), D2=(2-s,0);
pair X1=(s,2), X2=(2-s,2);
path reg=arc(P,2,-30,30)--C--arc(P,sqrt(6),45,135)--X2--arc(P,2,150,210)--A--arc(P,sqrt(2),225,315)--B--cycle;
fill(reg,rgb(1,0.93,0.75));
fill(P--X1--C--cycle,rgb(0.75,0.85,1)); fill(P--C2--X2--cycle,rgb(0.75,0.85,1));
fill(P--D2--A--cycle,rgb(0.75,0.85,1)); fill(P--A2--B--cycle,rgb(0.75,0.85,1));
draw(A--B--C--D--cycle,blue+linewidth(0.9));
draw(A2--B2--C2--D2--cycle,red+dashed+linewidth(0.9));
draw(reg,linewidth(1.3));
draw(P--B,gray); draw(P--X1,gray); draw(P--C,gray); draw(P--C2,gray); draw(P--X2,gray); draw(P--D2,gray); draw(P--A,gray); draw(P--A2,gray);
dot(P); dot(A); dot(B); dot(C); dot(D);
label("$P$",P,2*W); label("$A$",A,NE); label("$B$",B,SE); label("$C$",C,NE); label("$D$",D,NE);
label("$A'$",A2,S,red); label("$B'$",B2,N,red); label("$C'$",C2,NW,red); label("$D'$",D2,SW,red);
label("$\sqrt6$",P+1.25*dir(45),NW,fontsize(8pt)); label("$2$",P+1.5*dir(-30),N,fontsize(8pt));
label("$2$",P+1.5*dir(210),N,fontsize(8pt)); label("$\sqrt2$",P+0.75*dir(225),SE,fontsize(8pt));
\end{asy}

\textbf{The area.} The region is made of four sectors and four triangles with apex $P$ (shaded darker in the figure):
\begin{itemize}
\item sectors: $\dfrac12\cdot6\cdot\dfrac{\pi}{2}+\dfrac12\cdot2\cdot\dfrac{\pi}{2}+2\cdot\dfrac12\cdot4\cdot\dfrac{\pi}{3}=\dfrac{3\pi}{2}+\dfrac{\pi}{2}+\dfrac{4\pi}{3}=\dfrac{10\pi}{3}$;
\item the triangles with vertices $P,(s,2),C$ and $P,C',(2-s,2)$ (the image of $P,C,(2,s)$) each have a base of length $s-2=\sqrt3-1$ on a side at distance $\sqrt3$ from $P$: area $\dfrac{\sqrt3\left(\sqrt3-1\right)}{2}=\dfrac{3-\sqrt3}{2}$ each;
\item the triangles with vertices $P,D',A$ and $P,A',B$ each have a base of length $\sqrt3-1$ on the line $y=0$, at distance $1$ from $P$: area $\dfrac{\sqrt3-1}{2}$ each.
\end{itemize}
The total area is
\[\dfrac{10\pi}{3}+\left(3-\sqrt3\right)+\left(\sqrt3-1\right)=\dfrac{10\pi}{3}+2=\dfrac{1}{3}(10\pi+6),\]
with $\gcd(10,6,3)=1$. Hence $a+b+c=10+6+3=\boxed{\textbf{(C) }19}$.
